\documentclass[11pt]{scrartcl}
\usepackage{ebb}



\title {Humpty Dumpty}
\subtitle {Entrenamientos Nacionales Centros Pagmos}

\author{Emmanuel Buenrostro}

\begin{document}
\maketitle

\section{Teoria}
En construcción lmao
\begin{itemize}
    \item Humpty
    \item Las chorromil propiedades (incluye OMM 2010/5)
    \item Ejemplo 
    \item Dumpty
    \item Las propiedades 
    \item Ejemplo 
\end{itemize}

\section{Problemas }
\epigraph{I'm just a kid, what's your excuse?}{Steven Universe a Diamante Blanco}


\subsection{Nivel I}

\begin{problem}
    Considere el triángulo rectángulo $ABC$ con $\angle BAC = 90^\circ$. Sea $\ell$ la recta que pasa por $B$ y el punto medio del lado $AC$. Sea $\Gamma$ la circunferencia con diámetro $AB$. La recta $\ell$ y la circunferencia $\Gamma$ se intersectan en el punto $P$, distinto de $B$.

Demuestra que la circunferencia que pasa por los puntos $B$, $C$ y $P$ es tangente a la recta $AC$ en $C$.
\end{problem}

\begin{problem}
    Sea $ABC$ un triangulo acutangulo con $AB> AC$ y sea $O$ su circuncentro. Sea $P \neq O$ la intersección del circulo de diametro $AO$ con el circuncirculo de $BOC$. Demuestra que $AP$ es la simediana.
\end{problem}

\begin{problem} %[PAGMO 2021/2]
    
    Considere el triángulo rectángulo isósceles $ABC$ con $\angle BAC = 90^{\circ}$. Sea $\ell$ la recta que pasa por $B$ y el punto medio del lado $AC$. Sea $\Gamma$ la circunferencia con diámetro $AB$. La recta $\ell$ y la circunferencia $\Gamma$ se intersectan en el punto $P$, diferente de $B$. Muestre que la circunferencia que pasa por los puntos $A$, $C$ y $P$ es tangente a la recta $BC$ en $C$.
    
\end{problem}



\subsection{Nivel II}

\begin{problem} %[Ibero 2023/4]
Sean $B$ y $C$ dos puntos fijos en el plano. Para cada punto $A$ del plano, fuera de la recta $BC$, sea $G$ el baricentro del triángulo $ABC$. 

Determina el lugar geométrico de los puntos $A$ tales que
\[
\angle BAC + \angle BGC = 180^\circ.
\]

\textit{Nota: El lugar geométrico es el conjunto de todos los puntos del plano que satisfacen la propiedad.}
\end{problem}

\begin{problem}
    Sea $ABC$ un triángulo. La simediana desde $A$ corta nuevamente a la circunferencia circunscrita en $K$. Sea $K'$ la reflexión de $K$ respecto a $BC$. 

Demuestra que $AK'$ es una mediana.
\end{problem}

\begin{problem}
    En la configuración anterior, demuestra que el punto medio de $AK$ es el $A$-Dumpty.
\end{problem}

\begin{problem}
    Sea $P$ un punto variable sobre el lado $BC$ del triángulo $ABC$. Sean $M$ y $N$ puntos en $AB$ y $AC$, respectivamente, tales que $PM \parallel AC$ y $PN \parallel AB$.

Demuestra que, cuando $P$ varía sobre $BC$, la circunferencia circunscrita de $AMN$ pasa por un punto fijo.
\end{problem}
\subsection{Nivel III}

\begin{problem} %[USAMO 2008]
     Sea $ABC$ un triángulo acutángulo y escaleno. Sean $M$, $N$ y $P$ los puntos medios de $\overline{BC}$, $\overline{CA}$ y $\overline{AB}$, respectivamente.  

Las mediatrices de $\overline{AB}$ y $\overline{AC}$ cortan al rayo $AM$ en los puntos $D$ y $E$, respectivamente. Sean $BD$ y $CE$ rectas que se intersectan en un punto $F$ dentro del triángulo $ABC$.  

Demuestra que los puntos $A$, $N$, $F$ y $P$ son concíclicos.
\end{problem}

\begin{problem} %[ELMO 2014/5]
    Sea $ABC$ un triángulo con circuncentro $O$ y ortocentro $H$. Sean $\omega_1$ y $\omega_2$ las circunferencias circunscritas de los triángulos $BOC$ y $BHC$, respectivamente. 

Supón que la circunferencia con diámetro $\overline{AO}$ corta nuevamente a $\omega_1$ en $M$, y que la recta $AM$ corta otra vez a $\omega_1$ en $X$. De manera similar, supón que la circunferencia con diámetro $\overline{AH}$ corta nuevamente a $\omega_2$ en $N$, y que la recta $AN$ corta otra vez a $\omega_2$ en $Y$.

Demuestra que las rectas $MN$ y $XY$ son paralelas.
\end{problem}

\begin{problem} %[Ibero 2024/2]
    Sea $ABC$ un triángulo acutángulo y sean $M$ y $N$ los puntos medios de $AB$ y $AC$, respectivamente. Sea $D$ un punto en el interior del segmento $BC$ tal que $DB<DC$. Sean $P$ y $Q$ las intersecciones de $DM$ y $DN$ con $AC$ y $AB$, respectivamente. 

Sea $R \ne A$ la otra intersección de las circunferencias circunscritas de los triángulos $PAQ$ y $AMN$. Si $K$ es el punto medio de $AR$, demuestra que
\[
\angle MKN = 2\angle BAC.
\]
\end{problem}
\end{document}